\documentclass[11pt,a4paper]{article}

% --- Paquetes Fundamentales de Tipografía y Lenguaje ---
\usepackage[utf8]{inputenc}
\usepackage[T1]{fontenc}
\usepackage[spanish,es-nodecimaldot]{babel}
\usepackage{lmodern}
\usepackage{microtype}

% --- Geometría y Márgenes ---
\usepackage[top=2.5cm,bottom=2.5cm,left=2.5cm,right=2.5cm]{geometry}
\usepackage{fancyhdr}
\pagestyle{fancy}
\fancyhf{}
\fancyhead[L]{\small\textsc{GeoQC-Field} $\cdot$ Manual Técnico de Control de Calidad}
\fancyhead[R]{\small MASW y TRS $\cdot$ Dosier 2.0}
\fancyfoot[C]{\thepage}
\renewcommand{\headrulewidth}{0.4pt}
\renewcommand{\footrulewidth}{0.4pt}

% --- Paquetes Matemáticos y Símbolos ---
\usepackage{amsmath,amssymb,amsfonts,mathtools}

% --- Tablas Profesionales ---
\usepackage{booktabs}
\usepackage{tabularx}
\usepackage{multirow}
\usepackage{array}
\newcolumntype{L}[1]{>{\raggedright\arraybackslash}p{#1}}
\newcolumntype{C}[1]{>{\centering\arraybackslash}p{#1}}
\newcolumntype{R}[1]{>{\raggedleft\arraybackslash}p{#1}}

% --- Gráficos, Colores y Enlaces ---
\usepackage{graphicx}
\usepackage{xcolor}
\usepackage{tcolorbox}
\tcbuselibrary{skins,breakable}

\definecolor{brandblue}{RGB}{3, 105, 161}
\definecolor{branddark}{RGB}{15, 23, 42}
\definecolor{usablegreen}{RGB}{22, 128, 60}
\definecolor{reviewamber}{RGB}{183, 121, 0}
\definecolor{unusablegrey}{RGB}{100, 116, 139}
\definecolor{badred}{RGB}{185, 28, 28}

\usepackage{hyperref}
\hypersetup{
    colorlinks=true,
    linkcolor=brandblue,
    citecolor=brandblue,
    urlcolor=brandblue,
    pdftitle={Manual Técnico de Control de Calidad Geofísico de Campo: MASW y TRS},
    pdfauthor={GeoQC Development Team},
    pdfsubject={Control de Calidad Geofísico MASW y TRS}
}

% --- Cajas de Información Estilizadas ---
\newtcolorbox{principiobox}[1]{
    colback=brandblue!5!white,
    colframe=brandblue,
    fonttitle=\bfseries,
    title={#1},
    arc=3mm,
    boxrule=0.8pt,
    left=3mm, right=3mm, top=2mm, bottom=2mm
}

\newtcolorbox{warningbox}[1]{
    colback=reviewamber!8!white,
    colframe=reviewamber!90!black,
    fonttitle=\bfseries,
    title={#1},
    arc=2mm,
    boxrule=0.8pt,
    left=3mm, right=3mm, top=2mm, bottom=2mm
}

\newtcolorbox{limitbox}[1]{
    colback=unusablegrey!8!white,
    colframe=unusablegrey!90!black,
    fonttitle=\bfseries,
    title={#1},
    arc=2mm,
    boxrule=0.8pt,
    left=3mm, right=3mm, top=2mm, bottom=2mm
}

% --- Título del Documento ---
\title{
    \vspace{-1.5cm}
    \textbf{\Large SISTEMA DE CONTROL DE CALIDAD GEOFÍSICO EN CAMPO} \\[0.3cm]
    \textbf{\LARGE Manual Técnico de Procedimiento y Dosieres 2.0} \\[0.3cm]
    \large Análisis Multicanal de Ondas Superficiales (MASW) y Tomografía de Refracción Sísmica (TRS)
}
\author{\textbf{GeoQC Core Architecture} $\cdot$ Especificación Técnica de Procedimiento}
\date{Octubre 2026 $\cdot$ Revisión 2.0 (Desacoplamiento Diagnóstico)}

\begin{document}

\maketitle
\thispagestyle{empty}

\begin{abstract}
\noindent Este manual establece las directrices teóricas, matemáticas y operativas del sistema de control de calidad geofísico en campo (\textsc{GeoQC-Field}) para exploraciones sísmicas superficiales basadas en \textbf{MASW} y \textbf{TRS}. El protocolo responde al cambio metodológico formal del sistema: gran parte del control de calidad puede y debe estimarse de los propios registros de ondas, pero reconociendo de manera estricta los límites infranqueables de la interpretación física cuando faltan metadatos instrumentales. Se describe exhaustivamente: (1)~los insumos mínimos requeridos para cada técnica y la gestión de información faltante; (2)~el procedimiento algorítmico paso a paso, desde la integridad de trazas y análisis de ruido, hasta la extracción de crestas de dispersión, ambigüedad direccional en número de onda firmado $q \in [-q_N, q_N]$ y consistencia multidispáro; y (3)~la arquitectura de salida del \textbf{Dosier 2.0}, fundamentada en el desacoplamiento riguroso entre la \emph{Calidad diagnóstica por bandas} y la \emph{Suficiencia del objetivo ingenieril} ($V_{s30}$, profundidad o modelos de velocidad).
\end{abstract}

\vspace{0.3cm}
\tableofcontents
\newpage

% =========================================================================
\section{Introducción y Filosofía del Control de Calidad Diagnóstico}
% =========================================================================

\subsection{El Cambio de Paradigma: Del Bloqueo Normativo a la Calificación Interna}
Los sistemas tradicionales de control de calidad sísmico tienden a cometer uno de dos errores extremos:
\begin{enumerate}
    \item \textbf{Permisividad no calificada:} Declarar ``verde'' o ``aprobado'' cualquier registro del cual el software logre trazar una curva suave, sin verificar si el lóbulo principal del arreglo invade el límite espacial de Nyquist, si existen ondas en sentido inverso o si la curva cambia drásticamente al modificar la ventana de tiempo.
    \item \textbf{Bloqueo normativo excesivo:} Trasladar los requisitos finales de aceptación operativa (tales como verificación de cotas topográficas oficiales, conocimiento de los rails de saturación del convertidor ADC o modelos de inversión de $V_{s30}$) a la calificación diagnóstica primaria del sismograma en campo.
\end{enumerate}

El sistema \textsc{GeoQC-Field} adopta un enfoque equilibrado y teóricamente sólido: \textbf{no necesitamos conocer perfectamente el instrumento y el subsuelo para decidir si un registro aporta dispersión o primeros arribos interpretables, pero tampoco podemos recuperar información física faltante a partir de una simple matriz de amplitudes numéricas.}

\begin{principiobox}{Pregunta Central del Dosier Diagnóstico 2.0}
\centering
\emph{«¿Qué bandas o componentes de este registro son utilizables para extraer y comparar dispersión o tiempos de viaje, y qué tan estable es esa interpretación ante cambios de procesamiento y geometría?»}
\end{principiobox}

\subsection{Principio de las Dos Evaluaciones Desacopladas}
Todo informe o dosier emitido por el sistema genera dos conclusiones totalmente separadas:
\begin{itemize}
    \item \textbf{Dimensión 1: Calidad Diagnóstica por Banda / Traza.} Clasifica cada bin de frecuencia o canal en:
    \begin{itemize}
        \item \textcolor{usablegreen}{\textbf{UTILIZABLE*}}: Pasa los controles internos objetivos de coherencia, contraste cresta--fondo, ausencia de competidores espaciales y estabilidad frente a variantes de procesamiento.
        \item \textcolor{reviewamber}{\textbf{REVISAR}}: Presenta advertencias objetivas (lóbulo en Nyquist, picos competidores espaciales, sensibilidad al retirar canales o dirección no resuelta).
        \item \textcolor{badred}{\textbf{NO\_UTILIZABLE / SIN\_SOPORTE}}: Amplitud espectral nula o menos de dos canales activos.
        \item \textcolor{unusablegrey}{\textbf{NO\_EVALUABLE}}: Arreglos irregulares o geometrías incompletas que impiden aplicar el kernel analítico uniforme.
    \end{itemize}
    \item \textbf{Dimensión 2: Suficiencia del Objetivo.} Informa si el conjunto de datos demuestra la cobertura, profundidad o precisión requerida para un fin específico (e.g., $V_{s30}$ según NEHRP/Eurocódigo o profundidad de roca basal). Permanece como \textbf{PENDIENTE} (\texttt{NOT\_DEMONSTRATED}) mientras no se aporten las evidencias externas requeridas (plan de posiciones, identificación modal verificada, etc.).
\end{itemize}

Bajo esta separación, la ausencia de los voltajes de saturación ADC en la cabecera dejará abierta la sospecha de clipping sin confirmación instrumental, pero \textbf{no bloqueará} la extracción ni la calificación de estabilidad de una cresta de dispersión aparente.

\newpage
% =========================================================================
\section{Qué se Recibe para Cada Técnica: Insumos y Metadatos}
% =========================================================================

\subsection{Tabla Comparativa de Insumos}

\begin{table}[h!]
\centering
\small
\renewcommand{\arraystretch}{1.25}
\begin{tabularx}{\textwidth}{L{3.2cm} X X}
\toprule
\textbf{Elemento} & \textbf{MASW (Ondas Superficiales)} & \textbf{TRS (Refracción Sísmica)} \\
\midrule
\textbf{Formato nativo} & SEG-2, Seismic Unix (SU), texto tabular o binario con cabecera documentada. & SEG-2, SU, o texto ASCII multicanal con marcas temporales. \\
\addlinespace
\textbf{Datos de entrada obligatorios} & Matriz $N_{\text{canales}} \times N_{\text{muestras}}$ de amplitudes finitas y escalares de muestreo. & Matriz $N_{\text{canales}} \times N_{\text{muestras}}$ de amplitudes finitas y marcas temporales. \\
\addlinespace
\textbf{Tiempo obligatorio} & Intervalo de muestreo $dt$ (o tasa $f_s = 1/dt$). & $dt$, tiempo de pre-disparo ($T_{\text{pre}}$) y sincronización (time-break / canal trigger). \\
\addlinespace
\textbf{Geometría obligatoria} & Coordenadas relativas de receptores $x_r$ y posición de la fuente $x_s$ (para offset $x_1$). & Coordenadas $x_r, y_r, z_r$ de receptores y posiciones $x_s, y_s, z_s$ de fuentes directas e inversas. \\
\addlinespace
\textbf{Metadatos instrumentales deseables} & Frecuencia propia de geófonos (e.g. 4.5\,Hz), rango dinámico ADC, unidades físicas. & Frecuencia propia (e.g. 10--14\,Hz), impedancia de amortiguamiento, escala ADC. \\
\addlinespace
\textbf{Insumos de ruido} & Ventana pre-evento en el registro o archivo de ruido ambiente no excitado. & Ventana previa al time-break ($T_{\text{pre}} \ge 100$\,ms) o registros pasivos. \\
\addlinespace
\textbf{Objetivo ingenieril} & Profundidad requerida, $V_{s30}$ normativo o perfil de rigidez esperado. & Profundidad de basamento, modelo de 2 o más capas, objetivos de ripabilidad. \\
\bottomrule
\end{tabularx}
\caption{Comparativa de insumos mínimos y metadatos requeridos para MASW y TRS.}
\label{tab:insumos}
\end{table}

\subsection{Lo que NO se Puede Estimar de los Datos Numéricos}
\begin{limitbox}{Fronteras Infranqueables de la Ingesta Numérica}
Es un error común asumir que un algoritmo avanzado puede deducir cualquier parámetro faltante. Se establecen los siguientes límites:
\begin{enumerate}
    \item \textbf{Escala temporal y espacial:} Una matriz pura de números flotantes no contiene información intrínseca sobre si $dt = 1$\,ms o $dt = 0.5$\,ms, ni si la distancia entre geófonos es de 1\,m o 5\,m. Sin $dt$ y coordenadas explícitas, el registro es \textbf{inútil}.
    \item \textbf{Canales auxiliares y de disparo:} Si el canal 1 se utilizó como sensor de impacto (trigger mecánico) pero se mapea en las cabeceras como el primer geófono del arreglo, todo el análisis espacial se corromperá. La correspondencia física debe estar verificada en la libreta de campo.
    \item \textbf{Voltajes de saturación instrumental:} Sin los niveles máximos del conversor análogo-digital (ADC rails o escala completa en milivoltios), no se puede confirmar físicamente un clipping instrumental; sólo es posible levantar una \emph{sospecha de clipping} por deformación morfológica de la señal.
\end{enumerate}
\end{limitbox}

\newpage
% =========================================================================
\section{Procedimiento Metodológico Paso a Paso}
% =========================================================================

\subsection{Paso 1: Ingesta, Integridad Digital y Geometría}
El flujo inicia con la verificación estricta de la estructura digital del sismograma:
\begin{enumerate}
    \item \textbf{Finitud y dimensión:} Se comprueba que no existan valores \texttt{NaN}, \texttt{Inf} o muestras corruptas. El número de canales $N \ge 3$ y el número de muestras $M \ge 32$.
    \item \textbf{Mapeo geométrico:} Se calculan las posiciones ordenadas $x_1 < x_2 < \dots < x_N$, el espaciamiento inter-canal $dx_k = x_{k+1} - x_k$ y la apertura total del arreglo:
    \begin{equation}
        L = x_N - x_1
    \end{equation}
    \item \textbf{Evaluación de regularidad:} Se verifica si el arreglo es uniforme:
    \begin{equation}
        \max_k |dx_k - dx| \le \epsilon \quad (\text{con } \epsilon = 10^{-9}\,\text{m})
    \end{equation}
    Si el arreglo presenta huecos o espaciamientos irregulares, se conserva la coordenada real pero se inhabilita el uso directo de fórmulas de kernel uniforme, asignando la etiqueta \texttt{KERNEL\_UNSUPPORTED}.
\end{enumerate}

\subsection{Paso 2: Diagnóstico de Canal y Trazas Problemáticas}
Antes de cualquier transformada espectral, cada traza $u_i(t)$ ($i=1,\dots,N$) es sometida a controles individuales:
\begin{itemize}
    \item \textbf{Canales nulos o constantes:} Se verifica si $\max(u_i) - \min(u_i) = 0$ o si la traza está idénticamente en cero. Estos canales son marcados de inmediato como \texttt{MALA} e inhabilitados para el análisis espacial.
    \item \textbf{Potencia RMS nativa:} Se calcula la amplitud RMS centrada:
    \begin{equation}
        \text{RMS}_i = \sqrt{\frac{1}{M} \sum_{m=1}^{M} \left( u_i(t_m) - \bar{u}_i \right)^2 }
    \end{equation}
    Permite detectar anomalías de acoplamiento o canales hiperamplificados.
    \item \textbf{Detección morfológica de mesetas (Sospecha de Clipping):} Se detectan secuencias donde la primera derivada numérica se anula en amplitudes extremas:
    \begin{equation}
        u_i(t_m) = u_i(t_{m+1}) = \dots = u_i(t_{m+K}) = \pm \max_t |u_i(t)| \quad (K \ge 3)
    \end{equation}
    Al no disponer de los límites ADC del fabricante, el sistema reporta \textbf{sospecha de deformación por saturación}, sin pretender validar la calibración de fábrica.
\end{itemize}

\subsection{Paso 3: Análisis de Ruido y Coherencia de Fondo}
\begin{warningbox}{Distinción Fundamental: Contraste de Imagen Espectral vs. SNR Física}
No todo contraste de amplitud en el plano de dispersión equivale a una Relación Señal/Ruido (SNR) física. Usar la coda tardía del sismograma asumiéndola como ``ruido puro'' conduce a valores de SNR artificialmente optimistas y engañosos, pues la coda contiene reflexiones, retrodispersión y energía guiada tardía.
\end{warningbox}

El sistema adopta dos protocolos según la disponibilidad de datos:
\begin{enumerate}
    \item \textbf{Caso A: Existe ventana previa al evento ($T_{\text{pre}}$ representativo):}
    Se calcula la potencia cuadrática en la ventana de ruido $P_{\text{ruido}}$ y en la ventana de señal $P_{\text{señal}}$. La SNR condicional de la traza se define por:
    \begin{equation}
        \text{SNR}_i = 10 \log_{10} \left( \frac{P_{\text{señal}, i}}{P_{\text{ruido}, i}} \right)
    \end{equation}
    Se evalúa la estabilidad mediante la Desviación Absoluta Mediana (MAD) para evitar que picos transitorios contaminen la varianza.
    \item \textbf{Caso B: No existe ventana previa (MASW con registro disparado al inicio):}
    No se inventa un valor de SNR. En su lugar, el sistema estima el \textbf{contraste cresta--fondo} de la función de coherencia espacial $I(q)$ generada por el propio campo de ondas multicanal (descrito en el Paso 4A).
\end{enumerate}

\newpage
\subsection{Paso 4A: Extracción y Calidad Diagnóstica de Dispersión (MASW)}
Para cada frecuencia seleccionada $f$, se aplica la transformada de desfase (\emph{phase-shift}) en el dominio de número de onda firmado $q \in [-q_N, +q_N]$, donde $q_N = \frac{1}{2 dx}$ es el límite de Nyquist espacial:

\begin{enumerate}
    \item \textbf{Normalización espectral con piso relativo:}
    Sea $U(f, x_k)$ la transformada de Fourier de la señal centrada con ventana de Hann. Para evitar amplificar ceros numéricos o ruido asintótico, se define un piso relativo:
    \begin{equation}
        |U(f, x_k)| > \text{piso\_rel} \times \max_f |U(f, x_k)| \quad (\text{con piso\_rel} = 10^{-6})
    \end{equation}
    La fase normalizada se expresa como $\hat{U}(f, x_k) = \frac{U(f, x_k)}{|U(f, x_k)|}$.

    \item \textbf{Escaneo direccional de coherencia multicanal:}
    Para un número de onda firmado $q$, la coherencia espectral normalizada $I(f, q)$ se calcula como:
    \begin{equation}
        I(f, q) = \frac{1}{N} \left| \sum_{k=1}^{N} \hat{U}(f, x_k) \, e^{2\pi i \, q \, (x_k - x_1)} \right|, \quad q \in \left[ -\frac{1}{2dx}, +\frac{1}{2dx} \right]
    \end{equation}
    Nótese que la velocidad aparente en la dirección positiva de avance es $c = f / q$ (para $q > 0$). El escaneo en $q < 0$ evalúa la energía que viaja en sentido contrario hacia la fuente.

    \item \textbf{Resolución del arreglo y exclusión del lóbulo principal:}
    La función de respuesta de un arreglo lineal uniforme de $N$ receptores con espaciamiento $dx$ posee su primer cero analítico a una distancia de:
    \begin{equation}
        \Delta q_{\text{nulo}} = \frac{1}{N \cdot dx} = \frac{1}{L}
    \end{equation}
    Sea $q_0 = f / c_0$ la ubicación de la cresta seleccionada. Para no clasificar el hombro natural del lóbulo de difracción como un modo superior o un pico competidor, se define una \textbf{zona de exclusión estricta}:
    \begin{equation}
        \Omega_{\text{exclusión}} = \left\{ q \;\Big|\; |q - q_0| < \frac{1}{N \cdot dx} \right\}
    \end{equation}

    \item \textbf{Cálculo de atributos diagnósticos intrínsecos:}
    \begin{itemize}
        \item \textbf{Coherencia en cresta:} $I_{\text{cresta}} = I(f, q_0) \in [0, 1]$.
        \item \textbf{Fondo espectral:} Mediana de la coherencia fuera del lóbulo:
        \begin{equation}
            I_{\text{fondo}}(f) = \text{mediana}_{q \notin \Omega_{\text{exclusión}}} I(f, q)
        \end{equation}
        \item \textbf{Contraste cresta--fondo:}
        \begin{equation}
            C(f) = I_{\text{cresta}}(f) - I_{\text{fondo}}(f)
        \end{equation}
        \item \textbf{Pico competidor máximo (Rival Ratio):}
        Se buscan todos los máximos locales de $I(f, q)$ fuera de $\Omega_{\text{exclusión}}$ en todo el dominio firmado $[-q_N, q_N]$. Si existe un máximo rival en $q_{\text{rival}}$:
        \begin{equation}
            R(f) = \frac{I(f, q_{\text{rival}})}{I_{\text{cresta}}(f)}
        \end{equation}
    \end{itemize}

    \item \textbf{Controles de ambigüedad espacial y límites físicos:}
    \begin{itemize}
        \item Si $q_0 \le \frac{1}{N \cdot dx}$: La longitud de onda es comparable a la longitud total del tendido $L$. El lóbulo principal abarca $q=0$, por lo que la dirección de propagación no puede resolverse. Motivo: \texttt{DIRECTION\_UNRESOLVED}.
        \item Si $q_0 \ge q_N$: La cresta se encuentra en o fuera del límite de muestreo espacial. Motivo: \texttt{NYQUIST\_LIMIT}.
        \item Si $q_0 + \frac{1}{N \cdot dx} \ge q_N$: El lóbulo principal toca el límite de Nyquist, produciendo fuga de energía y distorsión. Motivo: \texttt{LOBE\_AT\_NYQUIST}.
    \end{itemize}
\end{enumerate}

\newpage
\subsection{Paso 4B: Primeros Arribos y Reciprocidad (TRS)}
Para la tomografía de refracción sísmica, la extracción de primeros arribos se fundamenta en:
\begin{enumerate}
    \item \textbf{Algoritmo STA/LTA con refinamiento AIC:}
    Se evalúa la razón entre el promedio de energía de ventana corta (STA) y ventana larga (LTA). En la vecindad del disparo, se aplica el Criterio de Información de Akaike (AIC) sobre la señal $u_i(t)$:
    \begin{equation}
        \text{AIC}(k) = k \log_{10}(\text{var}(u_i[1..k])) + (M - k - 1) \log_{10}(\text{var}(u_i[k+1..M]))
    \end{equation}
    El tiempo de llegada preliminar corresponde al mínimo global de $\text{AIC}(k)$.
    \item \textbf{Control estricto de reciprocidad de pares de disparos:}
    Si se dispone de un disparo directo en $A$ con receptor en $B$ ($t_{AB}$) y su disparo inverso en $B$ con receptor en $A$ ($t_{BA}$), la física del principio de Fermat exige reciprocidad cinemática:
    \begin{equation}
        |t_{AB} - t_{BA}| \le \tau_{\text{tol}} \quad (\text{típicamente } \tau_{\text{tol}} = 2\,dt \text{ a } 3\,dt)
    \end{equation}
    Las discrepancias mayores se reportan como error geométrico o deriva de reloj de disparo.
    \item \textbf{Monotonicidad aparente y filtros de velocidad física:}
    La curva dromocrónica o hodógrafa debe cumplir pendientes físicamente viables para ondas refractadas en el subsuelo ($V_P \in [300, 6500]$\,m/s). Pendientes negativas no justificadas por topografía son marcadas para revisión.
\end{enumerate}

\subsection{Paso 5: Estabilidad de la Curva ante Variantes de Procesamiento}
\begin{limitbox}{La Sensibilidad de Procesamiento NO es Desviación Estándar ($\sigma$)}
Modificar las opciones del software y comparar los resultados mide la \emph{sensibilidad del algoritmo frente a decisiones del operador}. No representa la dispersión estocástica debida al ruido del terreno ni la incertidumbre gaussiana del subsuelo.
\end{limitbox}

En MASW, el sistema evalúa la cresta base $c_0(f)$ (calculada con ventana de Hann y 24 canales) frente a tres variantes obligatorias:
\begin{enumerate}
    \item \textbf{Variante 1 ($c_{\text{rect}}$):} Transformada directa sobre ventana rectangular (sin atenuar los extremos del registro).
    \item \textbf{Variante 2 ($c_{\text{drop1}}$):} Exclusión del canal 1 (más cercano a la fuente), \textbf{conservando estrictamente las posiciones espaciales reales} de los canales 2 a $N$ (sin reindexar $x_2$ como si fuera el origen 0\,m).
    \item \textbf{Variante 3 ($c_{\text{unpen}}$):} Selección de la cresta por máximo absoluto puro por frecuencia, eliminando la penalización de curvatura de la programación dinámica.
\end{enumerate}
Se calcula la sensibilidad relativa máxima:
\begin{equation}
    S(f) = \max \left( \frac{|c_{\text{rect}}(f) - c_0(f)|}{c_0(f)}, \; \frac{|c_{\text{drop1}}(f) - c_0(f)|}{c_0(f)}, \; \frac{|c_{\text{unpen}}(f) - c_0(f)|}{c_0(f)} \right)
\end{equation}
Si $S(f) > 0.05$ (5\,\%), el punto se declara sensible al procesamiento (\texttt{PROCESSING\_SENSITIVE}).

\subsection{Paso 6: Consistencia Multidispáro entre Registros}
De acuerdo con \emph{Ólafsdóttir et al. (2018)}, evaluar la concordancia de curvas obtenidas con distintos offsets de fuente aporta respaldo metodológico sustancial, siempre que se reconozca que \textbf{distintos offsets no son réplicas idénticas del mismo experimento}:
\begin{itemize}
    \item Los offsets cortos ($x_1 = 5$--10\,m) iluminan mejor las altas frecuencias pero sufren sesgo de campo cercano en ondas largas.
    \item Los offsets largos ($x_1 = 20$--30\,m) atenúan los modos superiores y el campo cercano en bajas frecuencias, pero atenúan drásticamente las frecuencias altas por disipación intrínseca.
\end{itemize}
El sistema realiza la intersección de frecuencias comunes y computa la mediana y el valor máximo de discrepancia:
\begin{equation}
    \Delta c_{AB}(f) = |c_A(f) - c_B(f)|
\end{equation}
Estas diferencias se presentan como estadísticos observados, asignando el estado \texttt{uncertainty\_status: NOT\_CALIBRATED} para no generar una falsa ilusión de exactitud formal.

\newpage
% =========================================================================
\section{Dosier de Control de Calidad TRS (Refracción Sísmica)}
% =========================================================================

\subsection{Insumos Recibidos para TRS}
Para la técnica de Tomografía de Refracción Sísmica (TRS), el sistema recibe:
\begin{itemize}
    \item \textbf{Registros de sismogramas multicanal:} Archivos en formato estándar SEG-2 (\texttt{.sg2}) o SEG-Y, adquiridos con $N$ canales activos (habitualmente $N=24$) y múltiples posiciones de tiro distribuidas ($F_1, F_2, \dots, F_K$).
    \item \textbf{Geometría del tendido:} Coordenadas lineales de los receptores $x_r \in [0, L]$ con espaciamiento geofónico constante $dx$, y coordenadas de las fuentes $x_s$ (disparos directos, inversos e intermedios).
    \item \textbf{Sincronización y escala temporal:} Intervalo de muestreo $dt$ (típicamente $0.125$\,ms a $0.25$\,ms, $f_s \ge 4000$\,Hz) y marcas de disparo (\emph{time-break} o pre-disparo $T_{\text{pre}}$).
\end{itemize}

\subsection{Paso a Paso del Procesamiento Diagnóstico TRS}
El generador del Dosier TRS (\texttt{report\_trs\_dossier.py}) ejecuta las siguientes etapas analíticas:

\subsubsection{1. Evaluación y Calificación por Traza Individual}
Cada traza individual $u_{i,j}(t)$ (canal $j$ del disparo $i$) es analizada de manera autónoma:
\begin{enumerate}
    \item \textbf{Integridad de muestras:} Se verifica la ausencia de valores \texttt{NaN}, \texttt{Inf} o cadenas nulas. El canal debe cumplir $\text{valid\_fraction} = 1.0$.
    \item \textbf{Diagnóstico de sensor:} 
    \begin{itemize}
        \item Canal muerto: rango pico a pico $\text{PTP} = \max(u) - \min(u) = 0$.
        \item Canal constante: desviación estándar $\sigma < 10^{-12} \times \text{PTP}$.
    \end{itemize}
    \item \textbf{Amplitud y línea base:} Se remueve la componente continua (DC) y se calcula la amplitud RMS centrada nativa.
    \item \textbf{Sospecha de clipping morfológico:} Se detectan mesetas de $\ge 4$ muestras idénticas en extremos de amplitud. La fracción de clipping se califica bajo los umbrales locales de política (\texttt{TRACE\_QUALITY\_V1\_DRAFT}):
    \begin{equation}
        f_{\text{clip}} \le 0.001 \implies \text{BUENA}, \quad 0.001 < f_{\text{clip}} < 0.01 \implies \text{REGULAR}, \quad f_{\text{clip}} \ge 0.01 \implies \text{MALA}
    \end{equation}
    \item \textbf{Detección causal de primer arribo (P-Wave Onset):}
    Se implementa un detector causal STA/LTA (ventana corta de 2\,ms, ventana larga de 20\,ms) sobre la energía cuadrática $u^2(t)$. En la vecindad del disparo, se aplica el Criterio de Información de Akaike (AIC):
    \begin{equation}
        \text{AIC}(k) = k \log \left(\max(\text{var}(u[1..k]), \epsilon)\right) + (M-k) \log \left(\max(\text{var}(u[k+1..M]), \epsilon)\right)
    \end{equation}
    El arribo temporal $t_{\text{onset}}$ se obtiene del mínimo de $\text{AIC}(k)$, corrigiendo el retardo del pre-disparo.
    \item \textbf{Relación Señal/Ruido (SNR Neta en dB):}
    Se evalúa la potencia de ruido $P_{\text{ruido}}$ en la ventana previa al arribo ($[0, T_{\text{pre}}]$) y la potencia de señal $P_{\text{señal}}$ en la ventana posterior al arribo ($[T_{\text{pre}}, T_{\text{pre}} + 1500\,dt]$):
    \begin{equation}
        \text{SNR}_{\text{dB}} = 10 \log_{10} \left( \frac{P_{\text{señal}}}{P_{\text{ruido}} + \epsilon} \right)
    \end{equation}
    \item \textbf{Asignación de Calificación por Peor Hallazgo:}
    La función \texttt{geoqc.qc.quality.trace\_quality()} dictamina:
    \begin{itemize}
        \item \textcolor{usablegreen}{\textbf{BUENA}}: Integridad 100\%, $\text{SNR} \ge 12.0$\,dB, clipping $\le 0.1\%$, sin alertas de ruido.
        \item \textcolor{reviewamber}{\textbf{REGULAR}}: $\text{SNR} \in [6.0, 12.0)$\,dB, clipping moderado $\in (0.1\%, 1.0\%]$, o evidencia faltante.
        \item \textcolor{badred}{\textbf{MALA}}: Canal muerto, canal constante, $\text{SNR} < 6.0$\,dB, clipping $\ge 1.0\%$, o muestras no válidas.
        \item \textcolor{unusablegrey}{\textbf{NO\_EVALUABLE}}: Falta de soporte de muestras en la ventana de análisis.
    \end{itemize}
\end{enumerate}

\subsubsection{2. Evaluación y Calificación del Tendido Global}
Para calificar el tendido de refracción en su conjunto, el sistema aplica agregación rigurosa mediante \texttt{geoqc.qc.quality.assess\_quality()}:
\begin{enumerate}
    \item \textbf{Regla del denominador completo:} Todas las trazas físicas del tendido permanecen en el denominador, incluidas aquellas con evidencia incompleta o con arribo no resuelto. No se normaliza por subconjuntos ni se redondea antes del corte.
    \item \textbf{Intervalo de cota por evidencia pendiente:}
    Sea $N_{\text{total}}$ el número total de trazas, $N_{\text{buena}}$ las trazas calificadas como BUENA, y $N_{\text{pendiente}}$ aquellas con evidencia pendiente que no demostraron falla:
    \begin{equation}
        p_{\text{lower}} = \frac{N_{\text{buena}}}{N_{\text{total}}} \times 100\%, \quad p_{\text{upper}} = \frac{N_{\text{buena}} + N_{\text{pendiente}}}{N_{\text{total}}} \times 100\%
    \end{equation}
    \item \textbf{Decisión global del tendido:}
    \begin{itemize}
        \item \textcolor{usablegreen}{\textbf{BUENA}}: $p \ge 90.0\%$.
        \item \textcolor{reviewamber}{\textbf{REGULAR}}: $75.0\% \le p < 90.0\%$.
        \item \textcolor{badred}{\textbf{MALA}}: $p < 75.0\%$.
        \item \textcolor{unusablegrey}{\textbf{NO\_EVALUABLE}}: Si las cotas $p_{\text{lower}}$ y $p_{\text{upper}}$ caen en categorías distintas.
    \end{itemize}
    \item \textbf{Cobertura espacial de primeros arribos:}
    Se evalúa el porcentaje de canales del tendido que presentan arribo causal identificado ($N_{\text{arribos}} / N_{\text{total}}$).
    \item \textbf{Abanico de hodógrafas y ajuste de ramas de refracción:}
    Para cada disparo, se ajustan ramas lineales directas y refractadas mediante \texttt{fit\_branches()}:
    \begin{equation}
        t(x) = \begin{cases}
            \frac{x}{V_1}, & x < x_c \\
            \tau_2 + \frac{x}{V_2}, & x \ge x_c
        \end{cases}
    \end{equation}
    donde $x_c$ es la distancia de cruce y el espesor aparente del primer estrato es:
    \begin{equation}
        z_{\text{aparente}} = \frac{x_c}{2} \sqrt{\frac{V_2 - V_1}{V_2 + V_1}}
    \end{equation}
    \item \textbf{Validación de reciprocidad de Fermat:}
    Para pares conjugados de fuentes opuestas ($A \to B$ y $B \to A$), la función \texttt{reciprocal\_pair()} contrasta la discrepancia temporal $|t_{AB} - t_{BA}|$ contra la tolerancia operativa ($\le 3.0$\,ms).
    \item \textbf{Repetibilidad intra-fuente (Consistencia de impacto):}
    Se calcula la matriz de correlación de Pearson ($r$) entre impactos sucesivos apilados en cada estación de disparo para evaluar el efecto de compactación de placa y estabilidad mecánica.
    \item \textbf{Declaración de suficiencia del objetivo:}
    El dosier declara formalmente el estatus de suficiencia del objetivo como \textbf{PENDIENTE} (\texttt{NOT\_DEMONSTRATED}) para la inversión tomográfica 2D/3D mientras no se aporten ensayos independientes de pozo o comprobación de capas ocultas (ASTM D5777).
\end{enumerate}

\subsection{Estructura de Salida del Dosier TRS}
La ejecución de \texttt{report\_trs\_dossier.py} genera en la carpeta \texttt{evidence/trs/<SLUG>/}:
\begin{itemize}
    \item \textbf{\texttt{dossier\_trs.json}:} Archivo maestro JSON con digest canónico, hashes SHA-256 de archivos SEG-2, parámetros de protocolo, métricas de todas las trazas y decisiones de agregación global.
    \item \textbf{\texttt{dossier\_trs.pdf}:} Reporte gráfico multipágina en PDF con geometría del tendido, histograma de calidad, hodógrafas completas, matriz de calidad por canal vs. disparo, repetibilidad y reciprocidad de Fermat.
    \item \textbf{\texttt{dossier\_trs.html}:} Panel interactivo autónomo con estilos modernos en modo oscuro, visualizadores de gráficos, insignias de estado y tabla navegable de trazas.
    \item \textbf{\texttt{traces\_trs.csv}:} Registro tabular detallado de cada traza: canal, estación, offset, RMS, clipping, SNR neta, tiempo de arribo, calificación individual y motivo.
    \item \textbf{\texttt{picks\_trs.csv}:} Tiempos de primer arribo picados por canal y disparo.
    \item \textbf{\texttt{reciprocity\_trs.csv}:} Evaluación de pares de reciprocidad de Fermat con discrepancia $\Delta t$ y conformidad.
    \item \textbf{\texttt{repeatability\_trs.csv}:} Coeficientes de correlación mediana, mínima y máxima entre impactos consecutivos por estación.
    \item \textbf{\texttt{summary\_trs.csv}:} Síntesis de métricas del tendido, conteos de calidad, porcentajes y velocidades de rama.
\end{itemize}

\subsection{Resultados Obtenidos en Levantamientos TRS Reales}
Se procesaron tres levantamientos de campo de refracción sísmica:

\begin{table}[h!]
\centering
\small
\renewcommand{\arraystretch}{1.25}
\begin{tabularx}{\textwidth}{L{3.8cm} C{2.2cm} C{2.2cm} C{2.2cm}}
\toprule
\textbf{Parámetro de Diagnóstico TRS} & \textbf{KM 40+672} & \textbf{KM 44+552} & \textbf{KM 52+500} \\
\midrule
\textbf{Geometría del Tendido} & 24 ch $\times$ 5.0 m & 24 ch $\times$ 5.0 m & 24 ch $\times$ 2.5 m \\
\textbf{Longitud de Apertura ($L$)} & 115.0 m & 115.0 m & 57.5 m \\
\textbf{Estaciones de Disparo} & 9 fuentes ($F_1$ a $F_9$) & 9 fuentes ($F_1$ a $F_9$) & 11 fuentes ($F_1$ a $F_{11}$) \\
\textbf{Registros SEG-2 / Trazas} & 29 arch. / 696 trazas & 32 arch. / 768 trazas & 38 arch. / 912 trazas \\
\textbf{Finitud e Integridad} & 100\% finitas (0 NaN) & 100\% finitas (0 NaN) & 100\% finitas (0 NaN) \\
\textbf{Canales Muertos / Ceros} & 0 canales & 0 canales & 0 canales \\
\textbf{Trazas BUENAS} & 474 (68.1\%) & 480 (62.5\%) & \textbf{739 (81.0\%)} \\
\textbf{Trazas REGULARES} & 70 (10.1\%) & 114 (14.8\%) & 53 (5.8\%) \\
\textbf{Trazas MALAS} & 152 (21.8\%) & 174 (22.7\%) & 120 (13.1\%) \\
\textbf{Calificación Tendido Global} & \textcolor{badred}{\textbf{MALA ($<$75\%)}} & \textcolor{badred}{\textbf{MALA ($<$75\%)}} & \textcolor{reviewamber}{\textbf{REGULAR (81.0\%)}} \\
\textbf{Cobertura de Arribos} & \textbf{96.8\%} (674 onsets) & \textbf{99.3\%} (763 onsets) & \textbf{97.9\%} (893 onsets) \\
\textbf{Repetibilidad Media ($r$)} & 0.600 (picos $\ge$0.89) & 0.618 (picos $\ge$0.88) & \textbf{0.781 (Excelente)} \\
\textbf{Velocidad $V_1$ / $V_2$} & $\approx$811 / 2180 m/s & $\approx$850 / 2340 m/s & $\approx$980 / 2450 m/s \\
\textbf{Suficiencia del Objetivo} & \textbf{PENDIENTE} & \textbf{PENDIENTE} & \textbf{PENDIENTE} \\
\bottomrule
\end{tabularx}
\caption{Resultados consolidados del Dosier TRS sobre los levantamientos de campo reales.}
\label{tab:trs_km_results}
\end{table}

\newpage
% =========================================================================
\section{Dosier de Control de Calidad MASW (Ondas Superficiales a Detalle)}
% =========================================================================

\subsection{Insumos Recibidos para MASW}
Para el Análisis Multicanal de Ondas Superficiales (MASW), el sistema requiere:
\begin{itemize}
    \item \textbf{Registros multicanal con offset lateral declarado:} Sismogramas de impacto vertical registrados con un arreglo lineal de receptores donde la distancia fuente--primer geófono ($x_1$) es conocida y no nula (e.g. $x_1 = 5, 10, 20$\,m).
    \item \textbf{Espaciamiento y apertura:} Separación constante $dx$ y apertura $L = (N-1)dx$, calibrados para las velocidades de corte del terreno.
    \item \textbf{Parámetros de banda y velocidad:} Rango de frecuencias de interés ($f_{\min}$ a $f_{\max}$) y rejilla de velocidades candidatas ($c_{\min}$ a $c_{\max}$).
\end{itemize}

\subsection{Elementos a Detalle del Diagnóstico MASW}
El generador del Dosier MASW (\texttt{report\_masw\_quality.py}) evalúa minuciosamente cada uno de los componentes físicos y numéricos de la propagación:

\subsubsection{1. Transformada de Fase en Número de Onda con Signo}
Para evitar sesgos de asimetría direccional, la transformada de desfase (\emph{phase-shift}) se evalúa sobre el número de onda espacial con signo $q \in [-q_N, +q_N]$, donde $q_N = \frac{1}{2 dx}$ es el límite de Nyquist:
\begin{equation}
    I(f, q) = \frac{1}{N} \left| \sum_{k=1}^{N} \frac{U(f, x_k)}{|U(f, x_k)| + \epsilon} \, e^{2\pi i \, q \, (x_k - x_1)} \right|
\end{equation}
El cuadrante positivo ($q > 0$) corresponde a energía que se aleja de la fuente a velocidad $c = f/q$. El cuadrante negativo ($q < 0$) mapea energía retrodispersada o fuentes de ruido en el extremo opuesto del arreglo.

\subsubsection{2. Exclusión Geométrica del Lóbulo Principal Central}
El arreglo lineal de longitud $L = N \cdot dx$ posee una función de transferencia cuyo ancho de haz principal está delimitado por el primer cero analítico:
\begin{equation}
    \Delta q_{\text{nulo}} = \frac{1}{N \cdot dx} = \frac{1}{L}
\end{equation}
Para no confundir la caída natural del lóbulo con un segundo modo ni con ruido de fondo, se define la zona de exclusión alrededor de la cresta candidata $q_0$:
\begin{equation}
    \Omega_{\text{exclusión}} = \left\{ q \;\Big|\; |q - q_0| < \frac{1}{N \cdot dx} \right\}
\end{equation}

\subsubsection{3. Trazado de Cresta Aparente con Penalización de Continuidad}
La extracción de la curva de dispersión aparente fundamental $\hat{c}(f)$ se resuelve mediante programación dinámica minimizando una función de costo con penalización por salto de velocidad:
\begin{equation}
    J(c_1, \dots, c_K) = \sum_{k=1}^{K} \left[ 1 - I(f_k, c_k) \right] + \lambda \sum_{k=1}^{K-1} \left( \frac{c_{k+1} - c_k}{\Delta f \cdot s_0} \right)^2
\end{equation}
donde $\lambda = 0.02$ es la penalización de suavizado y $s_0 = 30$\,m/s/Hz es la escala de pendiente admisible.

\subsubsection{4. Atributos Diagnósticos Intrínsecos por Bin de Frecuencia}
En cada frecuencia analizada $f$, se calculan cuatro atributos objetivos:
\begin{itemize}
    \item \textbf{Coherencia en cresta:} $I_{\text{cresta}}(f) = I(f, \hat{c}(f)) \in [0, 1]$.
    \item \textbf{Contraste cresta--fondo:} Diferencia respecto a la mediana del espacio de lentitudes fuera de la zona de exclusión:
    \begin{equation}
        C(f) = I_{\text{cresta}}(f) - \text{mediana}_{q \notin \Omega_{\text{exclusión}}} I(f, q)
    \end{equation}
    \item \textbf{Razón de rival (Competitor Ratio):} Razón entre el segundo máximo espacial más alto y la cresta principal:
    \begin{equation}
        R(f) = \frac{\max_{q \notin \Omega_{\text{exclusión}}} I(f, q)}{I_{\text{cresta}}(f)}
    \end{equation}
    \item \textbf{Sensibilidad a variantes de procesamiento:} Variación relativa frente a tres variantes obligatorias (ventana rectangular, retiro del canal 1 preservando coordenadas y cresta libre no penalizada):
    \begin{equation}
        S(f) = \max_{v \in \{\text{rect}, \text{drop1}, \text{unpen}\}} \frac{|c_v(f) - \hat{c}(f)|}{\hat{c}(f)}
    \end{equation}
\end{itemize}

\subsubsection{5. Calificación Diagnóstica por Bandas}
Bajo la política \texttt{masw\_usability\_v1.json}, cada bin se clasifica objetivamente en:
\begin{itemize}
    \item \textcolor{usablegreen}{\textbf{UTILIZABLE*}}: Coherencia $\ge 0.60$, contraste $\ge 0.15$, razón de rival $\le 0.85$, sensibilidad $\le 5.0\%$, sin restricciones de Nyquist.
    \item \textcolor{reviewamber}{\textbf{REVISAR}}: Presencia de picos competidores ($R > 0.85$), contraste marginal ($C < 0.15$), o sensibilidad a variantes ($S > 5.0\%$).
    \item \textcolor{badred}{\textbf{NO\_UTILIZABLE}}: Cresta en el límite de Nyquist espacial ($q_0 \ge q_N$), lóbulo tocando Nyquist ($q_0 + 1/L \ge q_N$), lóbulo abarcando el origen $q=0$ (\texttt{DIRECTION\_UNRESOLVED}), o falta de amplitud espectral.
    \item \textcolor{unusablegrey}{\textbf{NO\_EVALUABLE}}: Variantes no computables o soporte geométrico irregular.
\end{itemize}

\subsubsection{6. Comparación y Consistencia Multidispáro}
Se evalúa la concordancia de las curvas de dispersión obtenidas a diferentes distancias de fuente ($x_1 = 10, 15, 20, 30$\,m), registrando la mediana de discrepancia $|\Delta c(f)|$ en los bins comunes.

\subsubsection{7. Declaración Formal de Suficiencia del Objetivo}
El dosier declara formalmente el estatus de suficiencia del objetivo como \textbf{PENDIENTE} (\texttt{NOT\_DEMONSTRATED}) para el perfil $V_S$ e índice $V_{s30}$, pues la transformación de velocidad de fase aparente a modelo estratificado exige comprobación de sitio (pozo o cross-hole) e identificación modal verificada.

\subsection{Estructura de Salida del Dosier MASW}
La ejecución de \texttt{report\_masw\_quality.py} genera en \texttt{evidence/masw\_quality/}:
\begin{itemize}
    \item \textbf{\texttt{dossier\_masw.json} / \texttt{dossier.json}:} Archivo JSON maestro con trazabilidad criptográfica SHA-256 de insumos, código y clasificaciones.
    \item \textbf{\texttt{dossier\_masw.pdf} / \texttt{dossier.pdf}:} Reporte formal multipágina con paneles de comparación multidispáro, mapa de calor de usabilidad por frecuencia, espectros $I(f, q)$ e $I(f, c)$ directos e inversos, y curvas de sensibilidad.
    \item \textbf{\texttt{dossier\_masw.html} / \texttt{dossier.html}:} Aplicación web autónoma en HTML5 con gráficos interactivos y tabla de intervalos de usabilidad.
    \item \textbf{\texttt{usability\_bands.csv}:} Bandas continuas de frecuencia clasificadas (\texttt{UTILIZABLE*}, \texttt{REVISAR}, \texttt{NO\_UTILIZABLE}).
    \item \textbf{\texttt{usability.csv}:} Registro detallado por bin: frecuencia, velocidad, coherencia, contraste, razón de rival, sensibilidad y motivos.
    \item \textbf{\texttt{comparisons.csv}:} Discrepancias $|\Delta c|$ entre pares de registros a diferentes offsets.
    \item \textbf{\texttt{points.csv}:} Coordenadas y velocidades de las crestas aparente en todas las variantes.
    \item \textbf{\texttt{traces.csv}:} Control de integridad y RMS de las trazas de entrada.
\end{itemize}

\subsection{Resultados Experimentales: Benchmark Oysand P1}
El procesamiento sobre el benchmark de Oysand ($dx = 2.0$\,m, 4 offsets de fuente de 10 a 30\,m) demostró:
\begin{itemize}
    \item \textbf{Bajas frecuencias ($<$8--10\,Hz):} Clasificadas en \texttt{NO\_UTILIZABLE} por \texttt{DIRECTION\_UNRESOLVED} ($q_0 \le 1/L$), dado que la longitud de onda excede la apertura del tendido.
    \item \textbf{Banda media ($\sim$10 a 22\,Hz):} Clasificada como \textcolor{usablegreen}{\textbf{UTILIZABLE*}} con coherencia $\ge 0.80$, contraste $\ge 0.35$ y discrepancia entre offsets inferior a 2\,m/s.
    \item \textbf{Muesca intermedia ($\sim$22--23\,Hz):} Clasificada en \textcolor{reviewamber}{\textbf{REVISAR}} por competencia de un modo superior incipiente (\texttt{COMPETING\_PEAK}) y sensibilidad al canal 1.
    \item \textbf{Altas frecuencias ($>$30\,Hz):} Clasificadas en \texttt{NO\_UTILIZABLE} por alcanzar el límite de Nyquist espacial (\texttt{NYQUIST\_LIMIT} y \texttt{LOBE\_AT\_NYQUIST}).
\end{itemize}

\newpage
% =========================================================================
\section{Guía Rápida para el Operador de Campo: TRS vs. MASW}
% =========================================================================

\begin{figure}[h!]
\centering
\small
\begin{tcolorbox}[colback=white,colframe=brandblue,arc=2mm,title=\textbf{Protocolo de Operación en Terreno $\cdot$ Decisión por Técnica}]
\begin{enumerate}
    \item \textbf{Para Ensayos TRS (Refracción):}
    \begin{itemize}
        \item \emph{Geometría recomendada:} Geófonos de 10--14\,Hz, espaciamiento $dx = 2.5$\,m a $5.0$\,m, apertura $L \ge 60$--120\,m.
        \item \emph{Disparos obligatorios:} Al menos 2 disparos en extremos conjugados ($F_{\text{inicio}}, F_{\text{fin}}$) para evaluar reciprocidad de Fermat, más 5 a 9 disparos intermedios para cruce de rayos.
        \item \emph{Control de golpe:} Realizar de 3 a 5 golpes por punto de tiro. El dosier TRS reportará la repetibilidad ($r \ge 0.70$ requerida).
        \item \emph{Dosier de entrega:} Ejecutar \texttt{report\_trs\_dossier.py --survey-dir <DIR>}. Verificar que la cobertura de arribos sea $\ge 90\%$ y que la reciprocidad cumpla $\Delta t \le 3$\,ms.
    \end{itemize}
    \item \textbf{Para Ensayos MASW (Ondas Superficiales):}
    \begin{itemize}
        \item \emph{Geometría recomendada:} Geófonos de 4.5\,Hz, espaciamiento $dx = 1.0$\,m a $2.0$\,m (nunca 5\,m si se busca caracterizar suelos blandos), apertura $L \ge 24$--48\,m.
        \item \emph{Disparos obligatorios:} Disparos laterales con al menos dos offsets diferentes (e.g. $x_1 = 10$\,m y $20$\,m) para evaluar consistencia multidispáro.
        \item \emph{Dosier de entrega:} Ejecutar \texttt{report\_masw\_quality.py}. Revisar en \texttt{dossier\_masw.html} las bandas en verde (\texttt{UTILIZABLE*}) y evitar picar en tramos con \texttt{LOBE\_AT\_NYQUIST}.
    \end{itemize}
\end{enumerate}
\end{tcolorbox}
\end{figure}

% =========================================================================
\section{Referencias Normativas y Bibliográficas}
% =========================================================================
\begin{enumerate}
    \item \textbf{Ólafsdóttir, E. A., Bessason, B., \& Erlingsson, S. (2018).} \emph{Combination of dispersion curves from multichannel analysis of surface waves (MASW).} Soil Dynamics and Earthquake Engineering, 113, 473--487. \href{https://doi.org/10.1016/j.soildyn.2018.05.025}{doi:10.1016/j.soildyn.2018.05.025}.
    \item \textbf{Park, C. B., Miller, R. D., \& Xia, J. (1999).} \emph{Multichannel analysis of surface waves.} Geophysics, 64(3), 800--808.
    \item \textbf{Foti, S., Lai, C. G., Rix, G. J., \& Strobbia, C. (2014).} \emph{Surface Wave Methods for Near-Surface Site Characterization.} CRC Press.
    \item \textbf{Aki, K., \& Richards, P. G. (2002).} \emph{Quantitative Seismology.} University Science Books.
    \item \textbf{Eurocódigo 8 / NEHRP (2020).} \emph{Provisions for Seismic Regulations for New Buildings and Other Structures ($V_{s30}$ Site Classification).}
\end{enumerate}

\end{document}
